\section*{الفصل \ref{ch:g9:linfunc} --- الدوال الخطية والتآلفية}

\begin{solution}{exo:g9:linfunc:1}
الخطية: $f$ و $k$ (على الصورة $ax$). والتآلفية غير الخطية: $g$
(\omterm{prop:g9:linfunc:affinegraph}{الميل} $3$، والترتيب عند المبدأ $-5$) و $m$ (ثابتة،
وميلها $0$). ولا هذه ولا تلك: $h$، بسبب المربّع.
\end{solution}

\begin{solution}{exo:g9:linfunc:2}
$f(0) = -3$؛ و $f(4) = 5$؛ و $f(-1) = -5$. و $f(x) = 9$ تعني
$2x - 3 = 9$، إذن $2x = 12$ و $x = 6$.
\end{solution}

\begin{solution}{exo:g9:linfunc:3}
$a = \frac{15}{6} = 2.5$، إذن $f(x) = 2.5x$. ومنه $f(10) = 25$
و $f(-2) = -5$.
\end{solution}

\begin{solution}{exo:g9:linfunc:4}
تمرّ $f$ بالمبدأ وميلها $2$؛ و $g$ لها \omterm{prop:g9:linfunc:affinegraph}{الميل} نفسه لكنّها
تقطع محور $y$ عند $3$؛ و $h$ متناقصة من $(0, 4)$ وميلها
$-1$. وتمثيلا $f$ و $g$ \emph{\omterm{def:g6:lines:perp}{مستقيمان متوازيان}} (\omterm{prop:g9:linfunc:affinegraph}{فالميل}
نفسه، والترتيبان عند المبدأ مختلفان).
\end{solution}

\begin{solution}{exo:g9:linfunc:5}
من أجل $f$: \omterm{prop:g9:linfunc:affinegraph}{الميل} $\frac{11 - 5}{3 - 1} = 3$، ثم
$5 = 3 \times 1 + b$ يعطي $b = 2$: أي $f(x) = 3x + 2$.

ومن أجل $g$: \omterm{prop:g9:linfunc:affinegraph}{الميل} $\frac{0 - 4}{2 - 0} = -2$، والقيمة $g(0) = 4$
هي الترتيب عند المبدأ أصلًا: أي $g(x) = -2x + 4$.
\end{solution}

\begin{solution}{exo:g9:linfunc:6}
زيادة $25\%$: $120 \times 1.25 = 150$.

وتخفيض $40\%$: $65 \times 0.6 = 39$.

والثمن الأصلي $p$ حيث $p \times 0.75 = 69$:
$p = \frac{69}{0.75} = 92$.
\end{solution}

\begin{solution}{exo:g9:linfunc:7}
\emph{1.} $f(x) = 0.05x + 10$.

\emph{2.} الاِشتراك B أرخص حين $25 < 0.05x + 10$، أي
$15 < 0.05x$، أي $x > 300$. فاِبتداءً من $301$ دقيقة ($5$ ساعات)
في الشهر، يفوز الاِشتراك B.
\end{solution}

\begin{solution}{exo:g9:linfunc:8}
\emph{1.} بعد ساعة: $2000 \times 1.1 = 2200$. وبعد ساعتين:
$2200 \times 1.1 = 2420$.

\emph{2.} تنطبق الزيادة الثانية على $2200$ لا على $2000$: فالنمو
بنسبة $10\%$ مرّتين يضرب في $1.1^2 = 1.21$، أي زيادة $21\%$، لا
$20\%$.
\end{solution}

\begin{solution}{exo:g9:linfunc:9}
\emph{1.} \omterm{prop:g9:linfunc:affinegraph}{الميل}: $\frac{1 - 3}{6 - 2} = \frac{-2}{4} = -0.5$:
فالدالة متناقصة.

\emph{2.} $f(x) = -0.5x + b$ مع $f(2) = 3$: $-1 + b = 3$، إذن $b = 4$
و $f(x) = -0.5x + 4$. والمخرج $0$: $-0.5x + 4 = 0$ يعطي $x = 8$.
\end{solution}

\begin{solution}{exo:g9:linfunc:10}
\emph{1.} يضرب التغيّران الثمن في
\[
\left(1 + \frac{t}{100}\right)\left(1 - \frac{t}{100}\right)
= 1 - \frac{t^2}{10000}
\]
(المتطابقة الثالثة مع $a = 1$ و $b = \frac{t}{100}$).

\emph{2.} من أجل $0 < t \leq 100$، يكون $\frac{t^2}{10000} > 0$،
فالعامل أصغر من $1$: أي أنّ الثمن النهائي أصغر من الأصلي. ومن أجل
$t = 10$: العامل $1 - \frac{100}{10000} = 0.99$، أي نقصان إجمالي
قدره $1\%$.
\end{solution}

\begin{solution}{pb:g9:linfunc:1}
\textbf{1.} \omterm{prop:g9:linfunc:affinegraph}{الميل}: $\frac{f(100) - f(0)}{100 - 0} =
\frac{212 - 32}{100} = 1.8$؛ والترتيب عند المبدأ $f(0) = 32$. ومنه
\[
f(C) = 1.8\,C + 32 .
\]

\textbf{2.} $f(20) = 68\,^\circ$F (اليوم المعتدل الشهير)؛
و $f(37) = 98.6\,^\circ$F؛ و $f(-10) = 14\,^\circ$F.

\textbf{3.} من $F = 1.8C + 32$:
$C = \frac{F - 32}{1.8}$. ومنه $50\,^\circ$F
$\to \frac{18}{1.8} = 10\,^\circ$C، و $98.6\,^\circ$F
$\to \frac{66.6}{1.8} = 37\,^\circ$C.

\textbf{4.} $f(0) = 32 \neq 0$، والقيمة $f(20) = 68$ ليست ضِعف
$f(10) = 50$: فليست خطية. والتمثيل البياني \omterm{def:g6:lines:objects}{مستقيم} لا يمرّ
بالمبدأ: أي أنّ $f$ \emph{تآلفية} لا خطية.

\textbf{5.} $1.8x + 32 = x$ يعطي $0.8x = -32$، إذن $x = -40$:
فعند أربعين تحت \omterm{def:g1:counting:zero}{الصفر}، يقرأ الميزانان $-40$ كلاهما. وعلى التمثيل
البياني، يقطع \omterm{def:g6:lines:objects}{مستقيم} $f$ \omterm{def:g6:lines:objects}{المستقيم} القطري $y = x$ هناك بالضبط.

\textbf{6.} $p_A(k) = 1.5k + 3$ و $p_B(k) = 2k$. ومن أجل $4$ كلم:
$9$ يورو مقابل $8$: فالسيارة B أرخص. ومن أجل $8$ كلم: $15$ مقابل
$16$: فالسيارة A أرخص.

\textbf{7.} $1.5k + 3 = 2k$ يعطي $k = 6$: فعند ستة كيلومترات
تكلّف السيارتان $12$ يورو. وبيانيًّا، يبدأ \omterm{def:g6:lines:objects}{مستقيم} B (المارّ
بالمبدأ، والأشدّ اِنحدارًا) تحت \omterm{def:g6:lines:objects}{مستقيم} A ويقطعه عند
$(6, 12)$: فتفوز B في الرحلات القصيرة، وتفوز A في الطويلة.

\textbf{8.} $N = 7 \times 20 - 30 = 110$ صريرًا في الدقيقة. ومن
$82 = 7T - 30$: $T = \frac{112}{7} = 16\,^\circ$C.

\textbf{9.} $v(t) = 600 - 15t$. والقيمة $240$:
$600 - 15t = 240$، أي $t = 24$ شهرًا. والقيمة $0$: أي $t = 40$
شهرًا. وبعد $40$ شهرًا تصير الصيغة سالبة --- لكنّ قيمة الهاتف
تقف عند \omterm{def:g1:counting:zero}{الصفر}: فلكل نموذج تآلفي نطاق يكون فيه ذا معنًى.

\textbf{10.} بالتركيب: $\times 1.25$ ثم $\times 0.8$ يضرب في
$1.25 \times 0.8 = 1$: فيعود الثمن إلى بدايته بالضبط. ولا معجزة:
فالنسب المئوية تعمل على مرجعين مختلفين (إذ ينطبق الخفض على الثمن
المرفوع)، وقد تصادف هنا أن تلاشى المضروبان --- وأمّا الحكاية
العامّة، بخسارتها المضمّنة $1 - \frac{t^2}{10\,000}$، فهي
\cref{exo:g9:linfunc:10}.

\textbf{11.} \omterm{prop:g9:linfunc:affinegraph}{الميل}: $\frac{17 - 20}{30} = -0.1$ سم في الدقيقة؛
و $h(t) = 20 - 0.1t$. وتموت عند $h(t) = 0$: أي $t = 200$ دقيقة
--- ثلاث ساعات وعشرون دقيقة من الضوء.

\textbf{12.} اِنطلق الثاني بتقدّم $10$ كلم (الترتيب عند المبدأ)؛
والأول أسرع ($24 > 18$ كلم/سا). واللحاق:
$24t = 10 + 18t$ يعطي $6t = 10$، أي $t = \frac53$ سا $= 1$ سا $40$
دق، عند الكيلومتر $24 \times \frac53 = 40$.

\textbf{13.} $f(x) + g(x) = (2 - 0.5)x + (3 + 7) = 1.5x + 10$:
فهي تآلفية، وميلها \emph{\omterm{def:g1:addition:def}{مجموع} الميلين}، وترتيبها عند
المبدأ \omterm{def:g1:addition:def}{مجموع} الترتيبين --- وهذا واضح بجمع الحدود المتشابهة في
$(mx + p) + (m'x + p')$.

\textbf{14.} إذا كان $m \neq 0$: فحلّ واحد بالضبط،
$x = -\frac pm$ --- إذ يقطع \omterm{def:g6:lines:objects}{المستقيم} غير الأفقي المحور مرّة
واحدة. وإذا كان $m = 0$ و $p \neq 0$: فلا حلّ --- إذ لا يلمس
المستقيمُ الأفقي فوق المحور أو تحته المحورَ أبدًا. وإذا كان
$m = 0$ و $p = 0$: فكل $x$ حلّ --- إذ يكون \omterm{def:g6:lines:objects}{المستقيم} \emph{هو}
المحور.

\textbf{15.} الثلاثية الأولى: المعدّلان $\frac{9 - 5}{3 - 1} = 2$
و $\frac{17 - 9}{7 - 3} = 2$: متساويان، فالنقاط على اِستقامة
واحدة (على $y = 2x + 3$). والثلاثية الثانية: $\frac{9-5}{2} = 2$
لكنّ $\frac{16 - 9}{3} = \frac73 \neq 2$: فليست على اِستقامة
واحدة. \omterm{def:g9:linfunc:affine}{والدوال التآلفية} هي بالضبط ذوات معدّل التغيّر الثابت ---
فعدد واحد، هو \omterm{prop:g9:linfunc:affinegraph}{الميل}، يروي الحكاية كلّها؛ أمّا المنحنيات
فيتغيّر معدّلها من نقطة إلى أخرى، ومطاردته تقود إلى المشتقّة، في
مجلّد الثانوية.
\end{solution}
